\section{研究启发与部署挑战}
\subsection{自监督与实战验证的融合}
讲者重申 ``Self-Feeding Chatbot'' 的工作流程：模型在对话后自己打标签，生成 verified response，再把高质量回答投入训练，从而实现 online self-improvement（SRT 774-884)。这个机制与 self-rewarding 训练的差别在于是真实用户交互的监督信号，而非纯粹在训练集中生成的 tasks。

\begin{warningbox}{部署中的鲁棒性风险}
Hallucination 仍然是 threat：如果 judge 自认为答案合理，而 hallucination 又恰好自洽，整个 self-improvement 就可能进入 ``强逻辑谬误'' 的放大环。需要引入第三方 evaluator 或多模态证据来锁定这类错误。
\end{warningbox}

\subsection{前沿问题与未来方向}
Weston 提到几个未来方向，包括 inference time compute for evaluation、interpretability、以及 ``self-aware'' 系统。具体来说，他说 ``Reason by actually interacting with people in the world, internet or itself''（SRT 3282-3284），并指出 self evaluation 的 compute 预算与 reasoning understand 之间的关系（SRT 3236-3252）。自知之明（self-aware）虽有些模糊，但至少意味着模型能在缺乏 judge 时主动寻求外部帮助。

\begin{importantbox}{未来工作三重点}
\begin{itemize}
    \item 推理 / 评估的 inference-time compute budget。
    \item 透明的 interpretation pipeline，让人类可以 trace reasoning。
    \item Self-aware 系统：模型在 uncertainty 高时知道去找 external agent。
\end{itemize}
\end{importantbox}

\subsection{本章小结}
研究必须兼顾 offline 的 self-training 与 online 的 verified feedback，并在 hallucination、compute 和 interpretability 之间寻求平衡。最终目标是打造既会推理又会自省的模型。

